\section{Related Work}
\label{sec:related-work}

\subsection{WHIR PCS Constructions}

WHIR~\cite{whir2024} was introduced by Hab\"ock, Levit, and Papini
as a Reed--Solomon proximity-testing IOP with super-fast
verification.  The Plonky3 implementation~\cite{plonky3} provides
the \texttt{p3-whir} adapter we build upon.  Our late-binding
extension is, to our knowledge, the first construction that allows
post-commit evaluation point derivation against a WHIR commitment
without modifying the upstream protocol.  The cryptographic
substrate (commitment writer, sumcheck rounds, STIR queries) is
identical to upstream WHIR; only the transcript-ordering wrapper
differs.

The VEIL paper~\cite{veil2026} provides a related construction:
a zero-knowledge multilinear PCS based on hash-based proximity
testing.  Lemma 4.7 of VEIL uses an interleaved-codeword commitment
and post-commit randomness in a similar spirit to our late-binding
adapter, though the motivation (zero-knowledge) and the protocol
shape differ.  We adopt the architectural insight ---
\emph{commit without opening points; sample batching randomness
post-commit} --- but restrict our scope to the non-zero-knowledge
case.

\subsection{Sumcheck-Based STARK Verification}

Sumcheck verification of AIR constraints (zerocheck) is a standard
technique in modern STARK constructions; we are not the first to
adopt it.  Our specific contribution is the integration into Ziren's
existing chip composition and the empirical validation that the
Phase 2a path produces verifiable proofs end-to-end.

The LogUp-GKR construction follows Hab\"ock~\cite{logup-gkr}.  Our
contribution is the discovery and remediation of a soundness bug
in our initial folding-only implementation: the combine identity
$(N \cdot D' + D \cdot N', D \cdot D')$ is degree-$2$ in each
coordinate, which only matches the verifier's degree-$1$
multilinear $N_{k+1}, D_{k+1}$ at Boolean points.  At random points
the prover can mount a forgery.  The fix is a per-layer sumcheck
reduction over a degree-$3$ round polynomial, restoring soundness
to $O(\log_2 N \cdot 3 / |\mathbb{F}_q|)$.

\subsection{SP1 Hypercube and Recursion Architectures}

SP1~\cite{sp1-hypercube} is a high-performance zkVM with a
hypercube-style architecture and an in-progress recursion-circuit
WHIR verifier (\texttt{crates/recursion/circuit/src/basefold/whir.rs},
$\sim 951$ lines, currently \texttt{\#![cfg(test)]}).  Our
recursion-circuit WHIR verifier scaffolding (\S\ref{sec:whir-late-binding})
follows SP1's protocol-shape closely: a $14$-helper inventory
maps onto SP1's per-component structure, with a clear port path
documented in our codebase.  The host-side cryptographic mechanisms
(zerocheck, LogUp-GKR, late-binding, cross-binding) are independent
of SP1's design and were derived per the Ziren architecture's
constraints.

\subsection{Brakedown and Spot-Check Protocols}

Our cross-binding (C.1) protocol is a Brakedown-style spot-check
argument~\cite{brakedown}: $K$ random codeword positions are sampled
post-commit, and the prover provides merkle paths against both
commits.  Brakedown originally targeted polynomial commitment and
R1CS proving directly; we use the same spot-check structure as a
\emph{cross-binding} layer over two independent commitments to the
same data.  Soundness is the standard $1 - (1 - 1/N)^K$.

\subsection{Jagged Polynomial Commitments}

The Jagged construction~\cite{jagged} packs variable-height matrices
into a single dense polynomial commit, amortising the per-matrix
overhead over a single PCS instance.  Our Phase 2c+ work combines
Jagged with our late-binding adapter via a sumcheck reduction that
ties per-matrix per-column row-MLE claims to a single dense MLE
claim $q(z^*)$, enabling per-shard commits at $57$--$85\times$
proof-size reduction over per-chip-per-column commits.  This sumcheck
reduction is, to our knowledge, novel; it adapts standard
sumcheck-on-product-of-multilinears to the per-chip-per-column
weighted-claim setting.
